\chapter{Propuesta de solucion}

\section{Vision general del sistema}

La solucion propuesta se organiza como un pipeline de procesamiento en tiempo real ejecutado localmente en el dispositivo movil. A partir de la imagen capturada por la camara, el sistema detecta landmarks de las manos, construye una representacion numerica de la secuencia gestual e infiere la seña correspondiente mediante un modelo liviano de clasificacion.

\section{Arquitectura propuesta}

La arquitectura se compone de los siguientes modulos principales:

\begin{itemize}
    \item Captura de imagen mediante CameraX.
    \item Deteccion de manos mediante MediaPipe Hand Landmarker.
    \item Extraccion de caracteristicas a partir de landmarks.
    \item Normalizacion de secuencias temporales.
    \item Clasificacion mediante TensorFlow Lite.
    \item Presentacion del resultado en texto y voz.
\end{itemize}

\section{Flujo de procesamiento}

Completar con una descripcion paso a paso del flujo: camara, frame, landmarks, vector de caracteristicas, buffer temporal, modelo, prediccion, texto y voz.

\section{Decisiones tecnologicas}

Completar justificando el uso de Android, Kotlin, CameraX, MediaPipe, TensorFlow Lite y Text-to-Speech.

\section{Justificacion de la ejecucion offline}

La ejecucion local permite reducir la dependencia de conectividad, mejorar la disponibilidad de la aplicacion en diferentes contextos y evitar el envio de imagenes o datos sensibles a servidores externos. Esta decision resulta especialmente relevante en aplicaciones vinculadas con accesibilidad y comunicacion personal.
